\documentclass{aediroum}

\title{Règlement de la vie étudiante}
\date{5 septembre 2025}

\begin{document}
\maketitle

\section{Préambule}\label{sec:preambule}

Afin d'assurer la pérennité des activités sociales de l'AÉDIROUM, le C.A. de l'AÉDIROUM a décidé de produire le présent règlement en avril 2016. Celui-ci vise à définir les activités du COUSSIN dans son rôle de support essentiel de la CVE de l'AÉDIROUM, ainsi qu'à définir les activités normalement offertes par celle-ci (CVE), sans lui donner l'obligation absolue de les organiser.

Il est à noter que les négociations ayant mené à la fusion des associations de premier cycle et de cycles supérieurs du Département d'Informatique et de Recherche Opérationnelle donnent l'obligation morale à l'AÉDIROUM d'offrir les beignes et café de cycles supérieurs (hebdomadairement) et les vins et fromages de l'AÉDIROUM (biannuellement).

Ce document se veut à la fois un document constitutif du comité permanent à la vie étudiante de l'AÉDIROUM, le COUSSIN, ainsi qu'un guide dans l'organisation de ses activités sociales.

L'objectif premier de ce document est de créer une structure facilitant l'organisation des diverses activités sociales de l'AÉDIROUM et de définir une répartition équitable et réaliste des tâches.

\section{Définitions}\label{sec:definitions}

\begin{description}
    \item[AÉDIROUM] Association Étudiante du Département d'Informatique et de Recherche Opérationnelle de l'Université de Montréal.
    \item[C.A.] Conseil d’administration de l’AÉDIROUM, tel que défini à l’article~3 de la Charte.
    \item[C.E.] Conseil exécutif de l'AÉDIROUM, tel que défini à l'article~6 de la Charte.
    \item[Charte] La Charte de l'AÉDIROUM.
    \item[COUSSIN] Comité Organisant l'Univers Social et Sportif INformatique.
    \item[Règlement] Règlement du guide du COUSSIN.
    \item[CVE] Coordination à la vie étudiante, tel que défini à l'article~6.1 de la Charte.
    \item[Assemblée générale d'élection] Assemblée générale des personnes membres de l'AÉDIROUM tenue au début du trimestre d'automne, en vertu de l'article 2 de la Charte.
    \item[Poste du COUSSIN] Un poste énuméré à l'article~\ref{sec:structure-coussin} du règlement.
    \item[Membre du COUSSIN] Une personne élue ou nommée à un poste du COUSSIN.
    \item[Activité] Une activité énumérée à l'article~\ref{sec:activites} du règlement.
\end{description}

\section{COUSSIN}\label{sec:structure-coussin}

\subsection{Fonctions du COUSSIN}\label{sec:coussin}

Le Comité Organisant l'Univers Social et Sportif INformatique (COUSSIN) est un comité permanent créé en vertu de l'article 7 de la Charte.

Le COUSSIN a pour but de veiller à la planification, l'organisation et la tenue des activités sociales de l'AÉDIROUM. Principalement il s'assure:

\begin{itemize}
    \item de favoriser les rencontres entre les membres de l'AÉDIROUM,
    \item de promouvoir un millieu social sain, inclusif, libre de harcèlement, libre de discrimination, et libre de favoritisme.
\end{itemize}

Il est composé de la CVE et des personnes adjointes à la CVE et des autres postes énumérés dans les sous-sections suivantes.

\subsection{Composition du COUSSIN}

\subsubsection{Coordination à la vie étudiante (CVE)}\label{sec:cve}

La personne coordonnatrice à la vie étudiante gère le COUSSIN et représente le COUSSIN auprès du C.E. de l’AÉDIROUM.
Elle s’assure d’obtenir les autorisations nécessaires à la tenue des événements, incluant ceux impliquant de l’alcool, et gère l’inventaire d’alcool conformément au Règlement de l’Université de Montréal et à la Loi sur les infractions en matière de boissons alcooliques.
Elle coordonne la planification des activités.
Elle s’assure de la bonne communication des activités en faisant des suivis auprès de la vice-présidence aux communications.

Les personnes adjointes à la coordination à la vie étudiante accompagnent la CVE à travers l'organisation et la planification des événements au fil de son mandat.
Elles voient au bon déroulement des activités qui sont placées sous leur supervision et participent avec la CVE aux différentes instances de la FAÉCUM conformément au Règlement sur l’affiliation à la FAÉCUM.

\subsubsection{Responsable des compétitions}\label{sec:responsable-competition}

La personne responsable des compétitions s'occupe de l'inscription, de la sélection ainsi que de la préparation nécessaire à la participation officielle de l'AÉDIROUM à diverses compétitions informatiques.

\subsubsection{Responsable des sports}\label{sec:responsable-sport}

La personne responsable des sports s'occupe d'organiser des activités sportives amicales ainsi que de la sélection et la participation officielle de l'AÉDIROUM à diverses activités sportives.

\subsubsection{Responsable des jeux}\label{sec:responsable-serveur}

La personne responsable des jeux s’occupe d’inventorier et de maintenir le matériel de jeu du local de l’AÉDIROUM, incluant les jeux de société, les consoles, les contrôleurs et les jeux pour les consoles.
Elle installe et tient à jour les jeux en ligne sur les serveurs informatiques de l’association. 
Elle organise des activités sociales autour des jeux.
Au besoin, elle recommande l’achat de nouveaux jeux.

\subsubsection{Responsable de la marchandise}

La personne responsable de la marchandise s'assure de la production d’une ligne de marchandise pour l'AÉDIROUM et de sa bonne distribution.
Elle est chargée de créer un visuel pour chaque nouvelle cohorte.
Elle gère l’inventaire de la marchandise ainsi que le site web de vente.
Elle peut former un comité pour l’assister dans ses tâches.

\subsubsection{Responsable des activités pour les cycles supérieurs}

La personne responsable des activités pour les cycles supérieurs organise différentes activités pour les cycles supérieurs telles que les beignes et café, les vins et fromages et les retraites de rédaction.
Elle s’assure également de la planification, de l’organisation et de la tenue des conférences MiDiro.

\subsubsection{Responsable des activités pour les personnes finissantes}

La personne responsable des activités pour les personnes finissantes est chargée de former un comité pour organiser les activités destinées aux personnes finissantes du baccalauréat.
Elle organise le bal annuel des personnes finissantes, en déléguant les différentes tâches à accomplir aux membres du comité et en faisant un suivi avec les membres, et veille au bon déroulement de l’événement.
Elle contacte les photographes et organise des séances de photographie pour les personnes finissantes et s’assure de la production, de la réception et de l’accrochage des mosaïques de finissants.

\subsection{Cumul de postes}\label{sec:cumul}

Les postes du COUSSIN sont cumulables.
En particulier, les personnes adjointes à la coordination à la vie étudiante peuvent occuper des postes du COUSSIN.

\subsection{Mandat}\label{sec:mandat-coussin}

Le mandat d'une personne membre du COUSSIN débute à son élection, ou à sa confirmation, jusqu'à l'assemblée générale d'élection subséquente.

\subsection{Élection}\label{sec:election-coussin}

Lors de l'assemblée générale d'élection, un vote sera tenu afin d'élire une personne membre de l'AÉDIROUM à chaque poste du COUSSIN.

\subsection{Nomination}\label{sec:nomination-coussin}

Advenant que des postes du COUSSIN autres que ceux tenus par des personnes membres du C.E. soient vacants après l'assemblée générale d'élection, la CVE sera habilitée à nommer un membre du COUSSIN, après consultation avec les membres restants du COUSSIN.
Cette nomination est assujettie à la confirmation du C.E. lors de sa prochaine séance.
Une personne membre nommée au COUSSIN qui n'a jamais été infirmée par le C.E. sera réputée en fonction jusqu'à sa confirmation.
Le C.E. ne peut nommer une personne membre du COUSSIN que si le poste de CVE est vacant.

\subsection{Destitution et démission}\label{sec:destitution-coussin}

Peut être destituée par résolution du C.E., toute personne membre du COUSSIN, autre que celles siégeant au C.E., qui ne remplit plus ses responsabilités.
Un membre du COUSSIN peut démissionner en remettant par écrit sa démission à la CVE ou à toute personne membre du C.A. de l'AÉDIROUM.

\section{Activités}\label{sec:activites}

\subsection{Ajout d'une activité}\label{sec:ajout-activite}

Une activité ajoutée au guide doit remplir un objectif, définir une fréquence et une personne responsable suggérée.

\subsection{Les vins et fromages des cycles supérieurs}\label{sec:vins-et-fromages}

\begin{description}
    \item[Description] Soirée vins et fromages au salon Maurice-Labbé. Les billets sont en vente plusieurs jours avant l'évènement et une tournée des laboratoires peut être faite pour faire la promotion de l'événement.
    \item[Objectif de l'activité] Favoriser des rencontres entre les membres des cycles supérieurs. Offrir un moment de socialisation et de réseautage.
    \item[Responsable suggéré] La personne responsable des activités pour les cycles supérieurs.
    \item[Fréquence] Une fois par session, en fonction du permis de réunion.
    \item[Note] La tenue des vins et fromages et des beignes et cafés furent une condition nécessaire à la création de l'AÉDIROUM (fusion des cycles supérieurs et du 1\textsuperscript{er} cycle).
\end{description}

\subsection{Les beignes et cafés des cycles supérieurs}\label{sec:beignes-et-cafes}

\begin{description}
    \item[Description] Petit intermède en milieu de semaine où les personnes étudiantes et professeures se retrouvent autour de beignes et de café.
    \item[Objectifs de l'activité] Favoriser des rencontres entre les personnes membres des cycles supérieurs. Offrir un moment de socialisation et de réseautage.
    \item[Responsable suggéré] La personne responsable des activités pour les cycles supérieurs.
    \item[Fréquence] Une fois par semaine, habituellement le mercredi 14h30 au salon Maurice-Labbé.
    \item[Note] La tenue des vins et fromages et des beignes et cafés furent une condition nécessaire à la création de l'AÉDIROUM (fusion des cycles supérieurs et du 1\textsuperscript{er} cycle).
\end{description}

\subsection{Les activités d’accueil}\label{sec:activites-accueil}

\begin{description}
    \item[Description] Journée où se côtoient plusieurs activités sociales, visant à créer des liens entre les participants par des activités loufoques et compétitives. Le tout est encadré par la FAÉCUM pour les réservations de terrain et de permis de réunion.
    \item[Objectifs de l'activité]
    \begin{itemize}
        \item[]
        \item Favoriser des rencontres entre les nouvelles personnes membres au baccalauréat
        \item Offrir un moment de socialisation
        \item Permettre aux nouvelles personnes étudiantes de rencontrer les anciens
        \item Permetter aux nouvelles personnes étudiantes de se familiariser avec le département et ses ressources
        \item Faire la promotion de l'AÉDIROUM auprès des nouvelles personnes étudiantes
        \item Mobiliser les nouvelles personnes étudiantes
    \end{itemize}
    \item[Responsables suggérés] La CVE et ses personnes adjointes.
    \item[Fréquence] Une fois par an. Tout juste avant la rentrée des nouvelles personnes étudiantes, c'est-à-dire fin août, début septembre.
\end{description}

\subsection{Les CS Games}\label{sec:cs-games}

\begin{description}
    \item[Description] Compétition informatique se déroulant pendant une fin de semaine complète, cumulant plusieurs épreuves dans différents domaines de l'informatique. L'évènement est organisé annuellement dans une université d'accueil différente.
    \item[Objectif de l'activité] Représenter le département au cours d'une compétition à teneur informatique
    \item[Responsable suggéré] La personne responsable des compétitions.
    \item[Fréquence] Une fois par an, aux alentours du mois de mars
\end{description}

\subsection{Les laits et biscuits du 1\textsuperscript{er} cycle}\label{sec:laits-et-biscuits}

\begin{description}
    \item[Description] Une convocation aux personnes membres est envoyée le jour avant pour les convier à se rencontrer au local de l'association. On met à la disposition des personnes membres des verres, environ 2$\times$2L de lait et des boîtes de biscuits.
    \item[Objectif de l'activité] Faire découvrir le local de l'association à de nouvelles personnes étudiantes du DIRO. Offrir un moment pour que les personnes du 1\textsuperscript{er} cycle se rencontrent et discutent.
    \item[Responsable suggéré] La personne adjointe à la CVE qui étudie en première année du baccalauréat.
    \item[Fréquence] Hebdomadaire.
\end{description}

\subsection{Les chalets des 1\textsuperscript{er cycle}}

\begin{description}
    \item[Description] Sortie dans un chalet loué par l’association pour deux nuits. Le départ et le retour se font par covoiturage.
    \item[Objectif de l'activité] Favoriser la socialisation entre les personnes étudiantes de premier cycle.
    \item[Responsable suggérés] La CVE et ses personnes adjointes.
    \item[Fréquence] Une fois à l’automne et une fois à l’hiver.
\end{description}

\subsection{Le chalet des cycles supérieurs}

\begin{description}
    \item[Description] Retraite de rédaction et activité de réseautage réservée aux personnes étudiantes aux cycles supérieurs. Le départ et le retour se font par covoiturage.
    \item[Objectif de l'activité] Réduire l’isolement aux cycles supérieurs, offrir un environnement propice à la rédaction, favoriser la socialisation entre les cycles supérieurs.
    \item[Responsable suggérés] La personne responsable des activités pour les cycles supérieurs.
    \item[Fréquence] Une fois à l’été.
\end{description}

\subsection{La cabane à sucre}\label{sec:cabane-a-sucre}

\begin{description}
    \item[Description] Sortie à la cabane à sucre. Le départ et le retour se font par covoiturage ou par autobus scolaire (1 ou 2).
    \item[Objectif de l'activité]
    \begin{itemize}
        \item[]
        \item Favoriser des rencontres entre les personnes membres. Offrir un moment de socialisation.
        \item Permettre aux personnes étudiantes internationales membres de l'AÉDIROUM de découvrir un aspect de la culture québécoise.
    \end{itemize}
    \item[Responsables suggérés] La CVE.
    \item[Fréquence] Une fois par an, au temps des sucres.
\end{description}

\subsection{Les MiDiro}\label{sec:midiros}

\begin{description}
    \item[Description] Présentation de sujet divers en informatique.

	Peut autant servir à la présentation de sujet de recherche qu'à la présentation d'outils de développement.

	Habituellement, une conférence d'une heure est donnée par une personne étudiante graduée qui vulgarise et présente son sujet de recherche à l’auditoire.
    \item[Objectif de l'activité]
    \begin{itemize}
        \item[]
        \item Être une plateforme pour présenter des sujets en informatiques hors des cours. Présenter le sujet de recherches de certaines personnes étudiantes graduées.
        \item Donner un aperçu aux personnes membres des projets de recherche du DIRO.
        \item Informer les personnes étudiantes du 1\textsuperscript{er} cycle sur différents aspects des études supérieures.
    \end{itemize}
    \item[Responsables suggérés] La personne responsable des activités pour les cycles supérieurs.
    \item[Fréquence] Une fois aux deux semaines ou en fonction des disponibilités des personnes étudiantes graduées.
\end{description}

\subsection{Le Carnaval de la FAÉCUM}\label{sec:carnaval}

\begin{description}
    \item[Description] Compétitions multiples s'échelonnant sur près d'un mois à la session d'hiver.
    \item[Objectif de l'activité] Compétitionner à l'échelle de l'université pour affronter d'autres programmes
    \item[Fréquence] Une fois par an, au mois de janvier et février.
    \item[Responsables suggérés] Les personnes adjointes à la CVE.
\end{description}

\end{document}
